\input{../tex-inputs/latex-korrekturansicht-vorspann}

\section[Theodor Herzl an Arthur Schnitzler, 6. 8. 1893]{L03832 Theodor Herzl an Arthur Schnitzler, 6. 8. 1893}
\nopagebreak\mylabel{L03832v}
\rehead{ }\normalsize\beginnumbering\briefempfaengerindex{Schnitzler, Arthur@\textsc{Schnitzler, Arthur}!zzzHerzl, Theodor@\emph{von Theodor Herzl}!1893-08-061@{6. 8. 1893}|(be}
\toendnotes[C]{\smallbreak\pagebreak[2]}
\correspDesc{Versand  durch Theodor Herzl am 6. 8. 1893 in Paris
\newline{}Erhalt  durch Arthur Schnitzler im Zeitraum [7. 8. 1893 – 11. 8. 1893?] in Wien}\toendnotes[C]{\smallbreak}
\Standort{CUL, Schnitzler, B 39.}
\physDesc{Brief, 1 Blatt, 1 Seite, 316 Zeichen
\newline{}Handschrift: schwarze Tinte, lateinische Kurrent
\newline{}Ordnung: mit Bleistift von unbekannter Hand nummeriert: »11« }
\buchAbdrucke{\weitereDrucke{Theodor Herzl: \emph{Briefe und
                        autobiographische Notizen 1866–1895}. Bearbeitet von Johannes Wachten in Zusammenarbeit mit Chaya Harel, Daisy Tycho und Manfred Winkler. Berlin, Frankfurt am Main, Wien: \emph{Propyläen} 1983, S. 538 (Briefe und Tagebücher. Herausgegeben von Alex Bein, Hermann Greive, Moshe Schaerf, Julius H. Schoeps und Johannes Wachten, 1).} }\toendnotes[C]{\smallbreak}
\pstart
           {\pb}\textcolor{gray}{\textbf{\textsc{\textcolor{brown}{Nouvelle Presse Libre}\orgindex{Neue Freie Presse@Neue Freie Presse|pw}{}\ledrightnote{\textcolor{brown}{Neue Freie Presse}}}}}\hfill \textcolor{pink}{\textcolor{gray}{\textbf{8, RUE DE MONCEAU}}}\oindex{8, rue de Monceau@\textbf{8, rue de Monceau}, \emph{Wohngebäude}|pw}{}\ledrightnote{\textcolor{pink}{8, rue de Monceau}}\pend
           
\pstart
           \textcolor{gray}{\textbf{D\textsuperscript{r}{ }TH. HERZL}}\pend
           
\pstart{}Lieber Freund!\pend\vspace{0.5em}
\pstart
           Die Wahrheit ist, dass ich unterwegs nach \textcolor{pink}{Oestreich}\oindex{Österreich@\textbf{Österreich}|pw}{}\ledrightnote{\textcolor{pink}{Österreich}} war und \label{K_L03832-1v}\edtext{von den
                  Ereignissen}{\lemma{\textnormal{\emph{von den
                  Ereignissen}}}\Cendnote{\textnormal{Am
                     2. 7. 1893 waren in \textcolor{pink}{Paris}\oindex{Paris@\textbf{Paris}, \emph{Hauptstadt}|pwk} Unruhen ausgebrochen, sodass \textcolor{blue}{Herzl}\pwindex{Herzl, Theodor 2.\,5.\,1860 Budapest – 3.\,7.\,1904 Edlach@\textsc{Herzl, Theodor} (2.\,5.\,1860 Budapest – 3.\,7.\,1904 Edlach), \emph{Schriftsteller, Journalist}|pwk} den Urlaub abbrechen und seinen Korrespondentenposten wieder
                  einnehmen musste.}}}\label{K_L03832-1} inmitten, in der schönen Mitten meines Urlaubs beim
               Kragen genommen u. nach \textcolor{pink}{Paris}\oindex{Paris@\textbf{Paris}, \emph{Hauptstadt}|pw}{}\ledrightnote{\textcolor{pink}{Paris}} zurückgeschleppt
               wurde.\pend
           
\pstart
           Wenn möglich komme ich im September nach \textcolor{pink}{Wien}\oindex{Wien@\textbf{Wien}, \emph{Verwaltungsgebiet}|pw}{}\ledrightnote{\textcolor{pink}{Wien}}. Dann sehen wir uns.\pend
           
\pstart
           In Eile aber immer herzlich ergeben{\\[\baselineskip]} Ihr{\\[\baselineskip]}\spacefill\mbox{Th Herzl}\pend
           \leftskip=0em{}
\pstart
           \noindent{}\textcolor{pink}{Paris}\oindex{Paris@\textbf{Paris}, \emph{Hauptstadt}|pw}{}\ledrightnote{\textcolor{pink}{Paris}}\pend
           
\pstart
           6 Aug. 93\pend
           \selectlanguage{ngerman}\endnumbering\briefempfaengerindex{Schnitzler, Arthur@\textsc{Schnitzler, Arthur}!zzzHerzl, Theodor@\emph{von Theodor Herzl}!1893-08-061@{6. 8. 1893}|)be}\mylabel{L03832h}  \newcommand{\dateiname}{L03832}\newcommand{\titel}{Theodor Herzl an Arthur Schnitzler, 6. 8. 1893}\newcommand{\editorInnen}{Selma Jahnke und Martin Anton Müller}\input{../tex-inputs/latex-korrekturansicht-abspann}
